\documentclass{article}

\usepackage[utf8]{inputenc}
\usepackage[T1]{fontenc}
\usepackage{amsmath,amssymb,amsthm}
\usepackage{mathtools}
\usepackage{hyperref}
\usepackage{booktabs}

\makeatletter
\newcommand{\citep}[2][]{%
  \begingroup
  \def\@tmpkey{#2}%
  \ifx\@tmpkey\@empty\else
    (\ifx\@empty#1\@empty\else #1, \fi
     \@nameuse{bibkey@\@tmpkey})%
  \fi
  \endgroup
}
\expandafter\def\csname bibkey@lasser2026micro\endcsname{Microfoundation}
\expandafter\def\csname bibkey@lasser2026exo\endcsname{Exogenous~Verification}
\expandafter\def\csname bibkey@lasser2026phase\endcsname{Phase~Redundancy}
\makeatother

\newtheorem{definition}{Definition}
\newtheorem{lemma}{Lemma}
\newtheorem{theorem}{Theorem}
\newtheorem{remark}{Remark}
\newtheorem{assumption}{Assumption}
\newtheorem{proposition}{Proposition}
\newtheorem{corollary}{Corollary}

\newcommand{\Bclip}{B_{\mathrm{clip}}}
\newcommand{\Kch}{K_{\mathrm{ch}}}
\newcommand{\Aadv}{\mathcal{A}_{\mathrm{adv}}}
\newcommand{\Pproxy}{P^{\mathrm{proxy}}}
\newcommand{\Ttruth}{T^{\mathrm{truth}}}
\newcommand{\epsnonresCoev}{\varepsilon_{\mathrm{coev}}^{\mathrm{nonres}}}
\newcommand{\Nev}{N_{\mathrm{ev}}}
\newcommand{\rhogap}{\rho_{\mathrm{gap}}}
\newcommand{\lammax}{\lambda_{\max}}
\newcommand{\taumeta}{\tau_{\mathrm{meta}}}
\newcommand{\Kchmulti}{K_{\mathrm{ch}}^{\mathrm{multi}}}
\newcommand{\bclk}{b_{\mathrm{clk}}}
\newcommand{\bcal}{b_{\mathrm{cal}}}

\title{C7 closed-form derivation, v3: \\
Lipschitz operator-norm formulation with separated rate cap}
\author{Paper 10 working draft for codex review}
\date{}

\begin{document}

\maketitle

\section{Goal and changes from v2}

\subsection{Goal}

Derive (C7) bounded co-evolution as a corollary of (C5.HOEFF), a
new explicit upper exposure-rate cap (C7.RATE), and the
verification protocol's channel-projection structure.  v3 makes
the intensivity claim cleaner than v2: it does \emph{not} require
the (C5.OVL) overlap-mass stability sub-clause that v2 introduced,
because the L1-normalized Lipschitz bound is automatically
intensive once $\Bclip$ and the rate cap are.

\subsection{Changes from v2 (codex review responses)}

\begin{enumerate}
\item \textbf{(C11.CLK) misuse fixed (P0).}  v2 used (C11.CLK) as
  an upper rate cap, but it is a lower throughput floor in the
  paper.  v3 introduces a separate (C7.RATE) condition for the
  upper rate cap.

\item \textbf{Lipschitz normalization fixed (P1).}  v2 used
  $\sup_{\|\Delta\|_1 \leq 1}$ which is not a proper operator norm
  unless the admissible class is star-shaped.  v3 uses
  $\sup_{\Delta \neq 0} |\Delta\mu_{j'}[\Delta]| / \|\Delta\|_1$,
  the operator norm of $\Delta \mapsto \Delta\mu_{j'}$.

\item \textbf{Perturbation model generalized (P1).}  v2's zero-sum
  $\Delta$ on $A_j$ couldn't represent Poisson rate shifts.  v3
  introduces an idle action $a_0$ representing
  ``no observable event,'' so rate changes are representable as
  signed redistributions including the idle bucket.

\item \textbf{Multi-label projection maps (P1).}  v2's binary
  incidence $I_j: A \to \{0,1\}$ was too coarse for ledger entries
  carrying multiple metadata tags.  v3 uses per-channel projection
  maps $\pi_j: A \to V_j$ where $V_j$ is channel $j$'s observable
  space; channel $j$ is engaged by action $a$ iff $\pi_j(a)$ is
  non-trivial.

\item \textbf{(C5.OVL) demoted to optional (P1).}  v2's
  overlap-mass stability sub-clause turned out to be unnecessary
  for the intensivity claim: $Q^*$ is automatically in $[0,1]$ as
  an L1 fraction.  v3 drops (C5.OVL) from the main proof and notes
  it is needed only for strict-smallness refinement (a tighter
  bound $\bar{M}^{\mathrm{cum}} \leq Q_{\max} \cdot \Bclip \cdot
  \lammax \cdot \taumeta$ instead of $\Bclip \cdot \lammax \cdot
  \taumeta$).

\item \textbf{Multi-channel aggregate bound added (P1).}  v2 used
  pairwise max, but Microfoundation's $K_{\mathrm{coev}}$
  aggregation may sum over source channels.  v3 states both
  pairwise and aggregate-incoming versions with explicit $\Kch -
  1$ factor.

\item \textbf{Pathwise window budget specified (P1).}  v2 left the
  ``unit perturbation maintained'' phrasing ambiguous.  v3 defines
  pathwise perturbation classes $\{\Delta_n\}_n$ with explicit
  budget conditions: per-step ($\sup_n \|\Delta_n\|_1 \leq 1$) or
  total-budget ($\sum_n \|\Delta_n\|_1 \leq 1$).

\item \textbf{Baseline projective consistency (P2).}  v2's
  baseline compatibility was underspecified for overlapping
  channels.  v3 adds a projective-consistency condition: $q_0$'s
  channel-$j$ projection $(\pi_j)_* q_0$ equals $p_0^{(j)}$.
\end{enumerate}

\section{Setup}

\subsection{Action space with idle bucket}

The verification ledger's per-step state space is $A = A^{\mathrm{ev}}
\sqcup \{a_0\}$, where $A^{\mathrm{ev}}$ is the set of recordable
event types and $a_0$ is the \emph{idle action} representing
``no observable event in this step.''  Per-step probability measures
on $A$ then represent both event distributions (under
$A^{\mathrm{ev}}$) and idle frequency (mass on $a_0$).  This
captures Poisson-rate-shift perturbations as redistributions
between $a_0$ and $A^{\mathrm{ev}}$.

\subsection{Channel projection structure}

The deployment monitors $\Kch$ channels indexed by $j \in
\{1, \ldots, \Kch\}$.  Each channel has:
\begin{itemize}
\item A \textbf{projection map} $\pi_j: A \to V_j \cup \{*_j\}$,
  where $V_j$ is channel $j$'s observable space (e.g., counts for
  Poisson, $\{0,1\}$ for Bernoulli, finite category set for
  multinomial) and $*_j$ is the ``trivial'' / null projection.  An
  action $a$ \emph{engages} channel $j$ iff $\pi_j(a) \in V_j$
  (equivalently, $\pi_j(a) \neq *_j$).
\item A baseline distribution $p_0^{(j)}$ on $V_j$ (the projected
  baseline; satisfies $p_0^{(j)} = (\pi_j)_* q_0$ by Assumption~2
  below).
\item A least-favorable adversarial alternative $p_1^{(j)}$ on
  $V_j$ at threshold shift $\eta_j$.
\item A clipped per-step LLR function
  $\ell^{(j)}: A \to [-\Bclip, \Bclip]$ that depends on $a$ only
  through $\pi_j(a)$, with $\ell^{(j)}(a) = 0$ whenever $\pi_j(a)
  = *_j$ (zero-outside-engagement convention).
\end{itemize}

The set of actions engaging channel $j$ is
$A_j := \{a \in A : \pi_j(a) \neq *_j\}$.  By the
zero-outside-engagement convention, $\ell^{(j)}(a) = 0$ for $a
\notin A_j$.

\subsection{Pairwise overlap}

For ordered channel pairs $(j, j')$ with $j \neq j'$, the
\emph{pairwise overlap set} is
\begin{equation}
  A_{j, j'} \;:=\; A_j \cap A_{j'}
  \;=\; \{a \in A : \pi_j(a) \neq *_j \text{ and } \pi_{j'}(a)
  \neq *_{j'}\}.
\end{equation}
Channels $j$ and $j'$ are \emph{action-disjoint} when $A_{j, j'} =
\emptyset$.  Multi-label entries (a ledger event carrying both
attestation and substrate metadata, say) appear naturally in
$A_{j, j'}$ if both projections are non-trivial.

\subsection{Per-step expected drift}

For an adversary $q$ on $A$, the per-step expected drift in channel
$j$ is
\begin{equation}
\label{eq:drift-def-v3}
  \mu_j(q) \;:=\; \mathbb{E}_{a \sim q}[\ell^{(j)}(a)]
  \;=\; \sum_{a \in A_j} q(a) \cdot \ell^{(j)}(a),
\end{equation}
the second equality by zero-outside-engagement.  Since
$\ell^{(j)}(a)$ depends on $a$ only through $\pi_j(a)$, this is
equivalently $\mathbb{E}_{v \sim (\pi_j)_* q}[\ell^{(j)}(v)]$.

\subsection{Baseline compatibility (with projective consistency)}

\begin{assumption}[Baseline ledger compatibility, projective form]
\label{ass:baseline-compat-v3}
There exists a baseline ledger distribution $q_0$ on $A$ such that
for every channel $j \in \{1, \ldots, \Kch\}$, the pushforward
under the channel projection map equals the per-channel baseline:
\begin{equation}
\label{eq:projective-baseline}
  (\pi_j)_* q_0 \;=\; p_0^{(j)}.
\end{equation}
The per-channel baselines $\{p_0^{(j)}\}_j$ must be
\emph{simultaneously realizable} as projections of a single $q_0$,
which constrains the per-channel baselines to agree on shared
projections (e.g., the marginal of $p_0^{(j)}$ on the shared
metadata between channels $j$ and $j'$ must equal the marginal
of $p_0^{(j')}$).  Verification protocols where per-channel
baselines are independently calibrated must check this consistency
at deployment activation.
\end{assumption}

\subsection{Admissible perturbations}

A signed measure $\Delta: A \to \mathbb{R}$ is \emph{admissible} if
$q_0 + \Delta \geq 0$ pointwise and $q_0 + \Delta$ is a probability
measure on $A$ (equivalently, $\sum_a \Delta(a) = 0$).  A signed
measure is \emph{$j$-supported} if $\Delta(a) = 0$ for $a \notin
A_j$; equivalently, the perturbation only changes mass on actions
engaging channel $j$ (potentially including the idle bucket if
$a_0 \in A_j$ via $\pi_j(a_0)$ being non-trivial; see below).

\begin{remark}[Idle-bucket convention]
\label{rem:idle}
Whether $a_0 \in A_j$ depends on the channel's protocol: for
``count rate'' channels (Poisson cooperative-output rate),
$\pi_j(a_0)$ is the ``no event in this step'' value, which is part
of $V_j$; so $a_0 \in A_j$.  For ``status'' channels (Bernoulli
attestation), $\pi_j(a_0) = *_j$ (no attestation event happened);
so $a_0 \notin A_j$.  This per-channel decision determines whether
rate-shift perturbations are representable in $\mathcal{D}_j$.
\end{remark}

The set of admissible $j$-perturbations is
\begin{equation}
  \mathcal{D}_j \;:=\; \{\Delta : A \to \mathbb{R}, \;\;\Delta
  \text{ is admissible and $j$-supported}\}.
\end{equation}

\section{Operational definition of \texorpdfstring{$M_{j \to j'}$}{M\_jj'}}

\subsection{Per-step coupling magnitude}

For $\Delta \in \mathcal{D}_j$, the induced shifts in per-channel
drift are
\begin{align}
  \Delta\mu_j[\Delta] &:= \mu_j(q_0 + \Delta) - \mu_j(q_0)
    = \sum_{a \in A_j} \Delta(a) \cdot \ell^{(j)}(a), \\
  \Delta\mu_{j'}[\Delta] &:= \mu_{j'}(q_0 + \Delta) - \mu_{j'}(q_0)
    = \sum_{a \in A_{j, j'}} \Delta(a) \cdot \ell^{(j')}(a),
\end{align}
the second equality because $\Delta$ is supported on $A_j$ and
$\ell^{(j')}$ vanishes off $A_{j'}$, so only $A_j \cap A_{j'} =
A_{j,j'}$ contributes.

\begin{definition}[Per-step coupling magnitude, operator-norm form]
\label{def:M-step-v3}
For $j \neq j'$, the \emph{per-step pairwise coupling magnitude}
of channel $j$ on channel $j'$ is the operator norm of the linear
map $\Delta \mapsto \Delta\mu_{j'}[\Delta]$ on the space of
admissible $j$-perturbations under the $\ell_1$ norm:
\begin{equation}
\label{eq:M-step-def-v3}
  M_{j \to j'}^{\mathrm{step}}(\mathcal{D})
  \;:=\; \sup_{\Delta \in \mathcal{D}_j,\; \Delta \neq 0}
  \;\frac{\big| \Delta\mu_{j'}[\Delta] \big|}{\|\Delta\|_1}.
\end{equation}

The \emph{per-step deployment coupling magnitude} is
$\bar{M}^{\mathrm{step}}(\mathcal{D}) := \max_{j \neq j'}
M_{j \to j'}^{\mathrm{step}}(\mathcal{D})$.
\end{definition}

\begin{remark}[Why this is a proper operator norm]
\label{rem:operator-norm}
The ratio $|\Delta\mu_{j'}[\Delta]| / \|\Delta\|_1$ is scale-
invariant under $\Delta \to t\Delta$ for $t \neq 0$ (both numerator
and denominator scale linearly).  The supremum over $\Delta \neq
0$ is therefore well-defined regardless of whether
$\mathcal{D}_j$ is star-shaped or scale-closed.  This is a proper
operator norm of the linear map $\Delta \mapsto \Delta\mu_{j'}[\Delta]$
restricted to the linear hull of $\mathcal{D}_j$.
\end{remark}

\section{Per-step closed-form bound}

\begin{lemma}[Per-step coupling bound]
\label{lem:per-step-v3}
Under (C5.HOEFF) and the zero-outside-engagement convention, for
every $j \neq j'$:
\begin{equation}
\label{eq:per-step-bound-v3}
  M_{j \to j'}^{\mathrm{step}}(\mathcal{D}) \;\leq\; \Bclip,
\end{equation}
unconditionally.  When channels are action-disjoint
($A_{j, j'} = \emptyset$), $M_{j \to j'}^{\mathrm{step}} = 0$.
\end{lemma}

\begin{proof}
For any $\Delta \in \mathcal{D}_j$ with $\Delta \neq 0$:
\begin{align*}
  \big| \Delta\mu_{j'}[\Delta] \big|
  &= \big| \sum_{a \in A_{j, j'}} \Delta(a) \ell^{(j')}(a) \big|
  \leq \sum_{a \in A_{j, j'}} |\Delta(a)| \cdot |\ell^{(j')}(a)|
  \\ &\leq \Bclip \sum_{a \in A_{j, j'}} |\Delta(a)|
  \leq \Bclip \cdot \|\Delta\|_1.
\end{align*}
Dividing by $\|\Delta\|_1 \neq 0$ and taking supremum gives
\eqref{eq:per-step-bound-v3}.  When $A_{j, j'} = \emptyset$, the
sum is empty and $\Delta\mu_{j'}[\Delta] = 0$ for all $\Delta$,
giving $M_{j \to j'}^{\mathrm{step}} = 0$.
\end{proof}

\begin{remark}[Tighter bound under (C5.OVL), if adopted]
\label{rem:tighter-v3}
If a deployment adopts the optional (C5.OVL) overlap-mass
stability condition (introduced in v2 but \emph{not} required for
intensivity), the supremum's shared-action concentration is
bounded:
\[
  \sup_{\Delta \in \mathcal{D}_j,\; \Delta \neq 0}
  \frac{\sum_{a \in A_{j, j'}} |\Delta(a)|}{\|\Delta\|_1}
  \;\leq\; Q^*_{j, j'}(\mathcal{D}) \leq Q_{\max} < 1,
\]
yielding $M_{j \to j'}^{\mathrm{step}} \leq Q_{\max} \cdot \Bclip$
in the strict-smallness regime.  This refinement matters when the
deployment claim requires a constant strictly less than $\Bclip$
(e.g., for stack-failure-probability budgeting); the intensivity
claim of Lemma~\ref{lem:per-step-v3} does not need it.
\end{remark}

\section{Cumulative cascade-window bound}

\subsection{Pathwise perturbation class}

A \emph{pathwise $j$-perturbation over the cascade window} is a
sequence $\{\Delta_n\}_{n=1}^{\Nev}$ where each
$\Delta_n \in \mathcal{D}_j$, applied at SPRT step $n$.  Two
budget conventions are operationally relevant:

\begin{itemize}
\item \textbf{Per-step budget:} $\sup_n \|\Delta_n\|_1 \leq 1$ ---
  the adversary can sustain a unit-magnitude perturbation each
  step.
\item \textbf{Total-budget:} $\sum_n \|\Delta_n\|_1 \leq 1$ --- the
  adversary has a fixed total perturbation budget across the
  window.
\end{itemize}

The cumulative cross-channel response is
\[
  \mu_{j'}^{\mathrm{cum}}\big[\{\Delta_n\}\big]
  \;:=\; \sum_{n=1}^{\Nev(\taumeta)} \Delta\mu_{j'}[\Delta_n].
\]

\subsection{New condition: upper exposure-rate cap (C7.RATE)}

\begin{assumption}[Upper exposure-rate cap, (C7.RATE)]
\label{ass:c7-rate}
There exists a deployment-class action rate cap $\lammax > 0$
independent of $|P|$, such that the SPRT exposure event count
satisfies
\begin{equation}
\label{eq:c7-rate}
  \Nev(\taumeta) \;\leq\; \lammax \cdot \taumeta
\end{equation}
deterministically (operator-enforced via rate limits / action
bounds).  This is distinct from (C11.CLK) which gives a lower
bound on $\Nev(\taumeta)$ for detection purposes.

Operationally, $\lammax$ is calibrated by Audit~7 from the
deployment's rate limits, action bounds, and channel-coupling
structure.
\end{assumption}

\subsection{Cumulative bound under per-step budget}

\begin{lemma}[Cumulative coupling bound, per-step budget]
\label{lem:cumulative-v3}
Under (C5.HOEFF), the zero-outside-engagement convention, and
(C7.RATE):

(a) \emph{Pairwise} cumulative coupling, per-step budget:
\begin{equation}
\label{eq:pairwise-cum-v3}
  \sup_{\{\Delta_n\}: \sup_n \|\Delta_n\|_1 \leq 1}
  \big| \mu_{j'}^{\mathrm{cum}}\big[\{\Delta_n\}\big] \big|
  \;\leq\; \Bclip \cdot \lammax \cdot \taumeta.
\end{equation}

(b) \emph{Aggregate-incoming} cumulative coupling, per-step
budget:
\begin{equation}
\label{eq:aggregate-cum-v3}
  \sup_{j', \{\{\Delta_n^{(j)}\}_n\}_j : \sup_n \sum_j
  \|\Delta_n^{(j)}\|_1 \leq 1}
  \big| \sum_{j \neq j'}
    \mu_{j'}^{\mathrm{cum}}\big[\{\Delta_n^{(j)}\}\big] \big|
  \;\leq\; (\Kch - 1) \cdot \Bclip \cdot \lammax \cdot \taumeta.
\end{equation}

The right-hand sides are deployment-class constants intensive in
$|P|$ over $\mathcal{D}$ ($\Bclip$ from (C5.HOEFF), $\lammax$ from
(C7.RATE), $\taumeta$ from \citep[Phase~Redundancy]{lasser2026phase}'s
scaling, $\Kch$ from (C5.MULT)).
\end{lemma}

\begin{proof}
\emph{Part (a).}  By Lemma~\ref{lem:per-step-v3}, $|\Delta\mu_{j'}[
\Delta_n]| \leq \Bclip \cdot \|\Delta_n\|_1 \leq \Bclip \cdot 1 =
\Bclip$ for each $n$ when $\sup_n \|\Delta_n\|_1 \leq 1$.  Summing:
\[
  \big| \mu_{j'}^{\mathrm{cum}} \big| \leq \sum_{n=1}^{\Nev(\taumeta)}
  |\Delta\mu_{j'}[\Delta_n]| \leq \Nev(\taumeta) \cdot \Bclip
  \leq \lammax \cdot \taumeta \cdot \Bclip,
\]
the last inequality by (C7.RATE).

\emph{Part (b).}  At each step $n$, the adversary's joint
perturbation across source channels has $\sum_j
\|\Delta_n^{(j)}\|_1 \leq 1$.  Each per-source-channel contribution
is bounded by $|\Delta\mu_{j'}[\Delta_n^{(j)}]| \leq \Bclip \cdot
\|\Delta_n^{(j)}\|_1$ via Lemma~\ref{lem:per-step-v3}.  Summing
over $j \neq j'$ at each step:
\[
  \big| \sum_{j \neq j'} \Delta\mu_{j'}[\Delta_n^{(j)}] \big|
  \leq \Bclip \cdot \sum_{j \neq j'} \|\Delta_n^{(j)}\|_1
  \leq \Bclip \cdot 1 = \Bclip.
\]
Summing over $n$ as in part (a):
\[
  \big| \sum_{j \neq j'} \mu_{j'}^{\mathrm{cum}}[\{\Delta_n^{(j)}\}] \big|
  \leq \Nev(\taumeta) \cdot \Bclip \leq \lammax \cdot \taumeta \cdot
  \Bclip.
\]
The factor $(\Kch - 1)$ enters only when the budget is
\emph{per-source-channel} ($\sup_n \|\Delta_n^{(j)}\|_1 \leq 1$ for
every $j$ rather than $\sum_j$): in that case each source channel
can saturate the unit budget, giving aggregate $\leq (\Kch - 1)
\cdot \Bclip \cdot \lammax \cdot \taumeta$.  We state both
budget conventions; which is operationally relevant depends on
the audit specification of admissible adversarial classes.
\end{proof}

\subsection{Connection to (C7) bounded co-evolution}

\begin{corollary}[(C7) bounded co-evolution as a corollary]
\label{cor:c7-corollary-v3}
Under (C5.HOEFF), the zero-outside-engagement convention, and
(C7.RATE), the deployment's cumulative coupling magnitude
$\bar{M}^{\mathrm{cum}}(\mathcal{D})$ over the cascade window
$\taumeta$ satisfies
\begin{equation}
  \bar{M}^{\mathrm{cum}}(\mathcal{D}) \;\leq\; C \cdot \Bclip \cdot
  \lammax \cdot \taumeta,
\end{equation}
where $C \in \{1, \Kch - 1\}$ depending on the budget convention
relevant to Microfoundation's $K_{\mathrm{coev}}$ aggregation.

In particular, $\bar{M}^{\mathrm{cum}}$ is bounded by a
deployment-class constant independent of $|P|$, which is the
content of Assumption~\ref{ass:bounded_coev_placeholder} (C7).
The corollary therefore promotes (C7) from a free-floating
deployment-class assumption to a derived consequence of
(C5.HOEFF), (C7.RATE), \citep[Phase~Redundancy]{lasser2026phase}'s
$\taumeta$ scaling, and (C5.MULT)'s channel partition.
\end{corollary}

(Where Assumption~\ref{ass:bounded_coev_placeholder} stands in for
Paper~10's existing Assumption labeled \texttt{ass:bounded\_coev}.)

\section{Discussion}

\subsection{What v3 delivers}

\paragraph{(C7) is now derived, not assumed.}  Under
(C5.HOEFF), (C7.RATE), and \citep[Phase~Redundancy]{lasser2026phase}'s
$\taumeta$ scaling, (C7)'s claim that $\bar{M}$ is bounded by a
$|P|$-independent constant follows from
Corollary~\ref{cor:c7-corollary-v3}.  The previous paper's
Assumption~\ref{ass:bounded_coev_placeholder} can be reframed as a
\emph{statement of (C7.RATE)} (operator-certified upper exposure-
rate cap) rather than a free-floating constant assumption.

\paragraph{Three primitive conditions, all already present or
inheritable.}  $\Bclip$ is from (C5.HOEFF).  $\taumeta$ is from
\citep[Phase~Redundancy]{lasser2026phase}.  $\Kch$ is from
(C5.MULT).  The new content is (C7.RATE), which adds an upper
exposure-rate cap separate from (C11.CLK)'s lower throughput floor.
This is the only new condition v3 introduces; (C5.OVL) from v2 is
not needed for the intensivity claim.

\paragraph{Audit~4 reinterpreted.}  Audit~4 currently measures
$\bar{M}$ empirically.  Under v3, Audit~4 can be reinterpreted as
verifying:
\begin{enumerate}
\item (C5.HOEFF)'s $\Bclip$ value (already part of Audit~7).
\item (C7.RATE)'s $\lammax$ value.
\item Channel projection structure: the $\pi_j$ maps are well-
  defined and respect the zero-outside-engagement convention.
\item (Optional) tighter (C5.OVL) ceiling $Q_{\max}$ for sharper
  bounds.
\end{enumerate}

\subsection{What v3 does not deliver}

\paragraph{Per-step vs.\ total-budget ambiguity persists.}
Lemma~\ref{lem:cumulative-v3} states bounds for two budget
conventions.  Which is operationally relevant depends on
Microfoundation's $K_{\mathrm{coev}}$ aggregation semantics, which
is not internal to this paper.  A future round should pin down
which budget convention $K_{\mathrm{coev}}$ uses.

\paragraph{Pairwise vs.\ aggregate-incoming ambiguity persists.}
Whether the deployment-relevant $\bar{M}^{\mathrm{cum}}$ is
pairwise (factor $C = 1$) or aggregate-incoming (factor $C = \Kch
- 1$) depends on how Microfoundation's composition correction
aggregates over source channels.

\paragraph{The strict-smallness refinement requires (C5.OVL).}  If
the deployment claim requires $\bar{M}^{\mathrm{cum}}$ strictly
less than the $\Bclip \cdot \lammax \cdot \taumeta$ ceiling
(e.g., for stack-failure-probability budgeting against $\beta +
\bclk + \bcal$ to leave room for $\bar{M}$-driven slack), the
optional (C5.OVL) condition $Q^*_{j,j'} \leq Q_{\max} < 1$ becomes
necessary.  v3 reports this as an optional refinement; the main
intensivity claim does not depend on it.

\subsection{Connection to (C11) capability-scaling caveat}

(C11)'s capability-scaling caveat noted that a more capable agent
may have a richer effective action space and find subtler
exploitation pathways.  v3's (C7.RATE) addresses this directly:
the upper exposure-rate cap $\lammax$ is operator-certified from
rate limits and action bounds, and any agent that exceeds
$\lammax$ falls outside the (C7.RATE) regime by definition.  The
capability-scaling caveat for $\rhogap$ is structurally parallel
and may admit a similar derivation: a closed-form bound on
$\rhogap$ from primitive operational parameters (rate limits +
clip radius + channel structure) would provide the same epistemic
shift for (C11) that v3 provides for (C7).

\section{Open questions for codex review}

\begin{enumerate}
\item \textbf{Is the operator-norm definition of
  $M_{j \to j'}^{\mathrm{step}}$
  (Definition~\ref{def:M-step-v3}) the right operational reading?}
  v3 uses $\sup_{\Delta \neq 0} |\Delta\mu_{j'}[\Delta]|/\|\Delta\|_1$.
  This is a proper operator norm, scale-invariant.  The L1
  normalization is inherited from the natural $\ell_1$ structure
  on the action space.  Is this consistent with what Audit~4
  empirically measures, or does Audit~4's perturbation methodology
  imply a different normalization?

\item \textbf{Is (C7.RATE) a clean formulation of the upper
  exposure-rate cap?}  v3 states it as a deterministic upper
  bound: $\Nev(\taumeta) \leq \lammax \cdot \taumeta$.  An
  alternative is a probabilistic bound (e.g., w.h.p.\ analogous to
  (C11.CLK)).  Which is cleaner for Microfoundation's composition
  semantics?

\item \textbf{Do the projection maps $\pi_j$ resolve the multi-label
  ledger entry concern from v2?}  v3 replaces binary incidence with
  per-channel projection maps.  Multi-label entries appear
  naturally in $A_{j, j'}$ via simultaneous non-trivial
  projections.  Are there ledger entry types that this still does
  not handle (e.g., conditional projections, hierarchical
  projections)?

\item \textbf{Does the projective consistency condition
  (Assumption~\ref{ass:baseline-compat-v3}) match what real
  verification protocols can guarantee?}  The condition requires
  $(\pi_j)_* q_0 = p_0^{(j)}$ for every $j$, with the per-channel
  baselines being simultaneously realizable.  In practice this
  constrains how per-channel baselines are calibrated.  Is this a
  reasonable practical constraint?

\item \textbf{Is the aggregate-incoming budget convention
  (per-source vs.\ total) a real operational distinction or a
  modeling artifact?}  v3 states both conventions and notes the
  distinction is downstream-resolved by $K_{\mathrm{coev}}$
  aggregation.  Is there a principled way to choose between them
  at the v3 level, or does the choice genuinely require a
  Microfoundation-internal answer?

\item \textbf{Are there v2-finding-leftovers that v3 missed?}
  Specifically: P0-1 (rate cap), P1-2 (Lipschitz norm), P1-3
  (Poisson rate-shift), P1-4 (pairwise sum), P1-5 (C5.OVL
  optional), P1-6 (multi-label projection), P1-7 (pathwise
  budget), P2-8 (baseline compatibility) should all be addressed
  in v3.  Is each adequately resolved?
\end{enumerate}

\section{Status}

v3 attempts to close the v2 P0/P1 gaps via:
\begin{itemize}
\item (C7.RATE) added as separate upper exposure-rate cap (P0-1).
\item Operator-norm Lipschitz definition (P1-2).
\item Idle-bucket convention for Poisson rate shifts (P1-3).
\item Aggregate-incoming bound with $(\Kch - 1)$ factor (P1-4).
\item (C5.OVL) demoted to optional sharper-bound condition (P1-5).
\item Per-channel projection maps replacing binary incidence (P1-6).
\item Pathwise perturbation class with per-step / total-budget
  conventions (P1-7).
\item Projective consistency in baseline compatibility (P2-8).
\end{itemize}

The structural shape: under (C5.HOEFF) clipping, (C7.RATE) upper
exposure-rate cap, and (C5.MULT) channel partition,
$\bar{M}^{\mathrm{cum}}$ is bounded by a $|P|$-independent
constant.  The closed form is $\Bclip \cdot \lammax \cdot \taumeta$
(pairwise) or $(\Kch - 1) \cdot \Bclip \cdot \lammax \cdot \taumeta$
(aggregate-incoming).

\end{document}
